\clearpage\SourceLeaf{146}%
\pdfbookmark[1]{Gauss an von Lindenau}{lindenau1808nov}
\fancyhead[L]{\small GAUSS AN VON LINDENAU}
\begin{center}
{\small GAUSS AN VON LINDENAU.\par}
\medskip\hrule\medskip
{\small Monatliche Correspondenz zur Beförderung der Erd- und Himmels-Kunde,\par
herausgegeben von Freiherrn \textsc{von Zach}. Januar 1809.\par}
\medskip\hrule
\end{center}\medskip
\begin{flushright}Göttingen, den 30. November 1808.\end{flushright}
Für die Beobachtungen von Ceres und Juno, die Sie mir so freundlich mitgeteilt haben, danke ich Ihnen herzlich. Auch für die Kometenbeobachtungen aus Marseille bin ich Ihnen dankbar, obwohl dabei \emph{sehr} große Fehler unterlaufen zu sein scheinen und sich deshalb schwerlich etwas daraus schließen lassen wird. Heute möchte ich Ihnen noch einige Bemerkungen zu der Aufgabe mitteilen, über die Sie meinen Aufsatz in das Oktoberheft der \emph{Monatl. Corresp.} aufgenommen haben. \emph{Reisenden} Beobachtern würde ich vor allem eine andere Aufgabe als \emph{diese} empfehlen. Ich habe sie am Anfang jenes Aufsatzes erwähnt und in einem vor Kurzem hier gedruckten Programm behandelt, von dem ich Ihnen gerne ein Exemplar beilege. Die in der \emph{Monatl. Corresp.} behandelte Aufgabe wird besonders für Beobachter brauchbar sein, die mit einem weniger vollkommenen Instrument ihren Beobachtungsort möglichst genau bestimmen und zugleich die Fehler des Instruments ermitteln wollen. In diesem Fall findet man ohne jede Schwierigkeit für jede Stunde viele geeignete Sterne, wenn man für die am leichtesten erkennbaren Sterne vorab eine Höhentabelle berechnet. Sie kann auch sonst nützlich sein und lässt sich sehr bequem als Karte gestalten. Man braucht sich keineswegs auf Sterne erster Größe zu beschränken.
\SourceLeaf{147}%
An meinem Beispiel sehen Sie, dass sich auch Sterne zweiter Größe verwenden lassen. Mit etwas Übung lassen sich bei fehlendem Mondschein sogar noch Sterne dritter und vierter Größe gut beobachten, obwohl man dann natürlich sorgfältiger sein muss, um \emph{Verwechslungen} zu vermeiden. Was meine Lösung selbst angeht, so bin ich sofort selbst darauf gekommen, als ich nach einem Verfahren suchte, das sich möglichst bequem in der Praxis anwenden lässt. Es fiel mir nicht ein, anderswo nach Lösungen einer Aufgabe zu suchen, mit der sich meines Wissens noch niemand beschäftigt hatte. Erst nachdem ich meinen Aufsatz abgeschickt hatte, bemerkte ich, dass eine andere Aufgabe zwar einen ganz anderen Zweck hat, sich aber im Wesentlichen genau gleich lösen lässt. Dabei sind aus drei \emph{heliozentrischen} Positionen eines Sonnenflecks die Lage des Sonnenäquators und zugleich die Deklination des Flecks zu bestimmen. Letztere entspricht in unserer Aufgabe der Höhe $h$. Bekanntlich haben sich viele Geometer mit dieser Aufgabe beschäftigt, doch hat niemand eine so elegante Lösung gegeben wie \textsc{Cagnoli}. \textsc{La Lande} führt sie in seiner \textit{Astronomie} an. Was die Eleganz betrifft, ziehe ich diese Lösung meiner eigenen vor. Ich glaube allerdings, dass sich die Lösung selbst eleganter aus den drei Grundgleichungen ableiten lässt, als \textsc{Cagnoli} es getan hat. Auch in der praktischen Handhabung stehen alle anderen von \textsc{La Lande} angeführten Lösungen weit hinter der von \textsc{Cagnoli} zurück. Meine Lösung wäre noch etwas bequemer, wenn man nur $\varphi$ und $k$ sucht, die von \textsc{Cagnoli} dagegen etwas kürzer, wenn man auch $h$ wissen will. Im ersten Fall brauche ich 18 Logarithmen und \textsc{Cagnoli} 21, im zweiten brauche ich 24 und \textsc{Cagnoli} wieder nur 21. Dabei ist allerdings zu beachten, dass bei \textsc{Cagnoli} alle 21 zu fast ebenso vielen, nämlich 19 \emph{verschiedenen} Winkeln gehören. Bei meiner Lösung schlägt man dagegen für viele Winkel Sinus und Kosinus oder andere trigonometrische Funktionen zugleich nach. (Man schlägt also an 14 verschiedenen Stellen der Tafeln nach, wenn man nur $\varphi$ und $k$ sucht, und an 18, wenn man auch $h$ wissen will.) Das macht durchaus einen Unterschied, besonders wenn man die ausgezeichneten Tafeln von \textsc{Taylor} verwendet, bei denen sich fast die ganze Arbeit auf das Nachschlagen beschränkt. Für unsere Aufgabe lauteten die Formeln nach \textsc{Cagnoli}s Lösung wie folgt.

Man berechne die drei Hilfswinkel $A,A',A''$ aus den Gleichungen
\SourceLeaf{148}%
\[
\begin{aligned}
\operatorname{tang}A&=\frac{\sin\frac12(\delta''-\delta')}{\cos\frac12(\delta''+\delta')}\operatorname{cotang}\tfrac12(t''-t')\\
\operatorname{tang}A'&=\frac{\sin\frac12(\delta-\delta'')}{\cos\frac12(\delta+\delta'')}\operatorname{cotang}\tfrac12(t-t'')\\
\operatorname{tang}A''&=\frac{\sin\frac12(\delta'-\delta)}{\cos\frac12(\delta'+\delta)}\operatorname{cotang}\tfrac12(t'-t)
\end{aligned}
\]
bestimme dann
\[
\operatorname{tang}B=\frac{\sin\frac12(\delta'-\delta)}{\cos\frac12(\delta'+\delta)}\operatorname{cotang}(A-A')
\]
dann gilt
\[
k=\tfrac12(t'+t)-B
\]
Setzt man ferner $A'+A''-A=C$ und nimmt
\[
\begin{aligned}
\operatorname{tang}D&=\frac{\cos\frac12(t-k+C)}{\cos\frac12(t-k-C)}\operatorname{tang}(45^0+\tfrac12\delta)\\
\operatorname{tang}E&=\frac{\sin\frac12(t-k-C)}{\sin\frac12(t-k+C)}\operatorname{cotg}(45^0+\tfrac12\delta)
\end{aligned}
\]
\expandafter\def\csname HistoricalNoteText6\endcsname{[Handschriftliche Bemerkung:]
\[
\begin{aligned}
\frac{\cos\varphi}{\cos h}\sin(t-k)&=\sin(A'+A''-A)\\
\frac{\cos\varphi}{\cos h}\sin(t'-k)&=\sin(A''+A-A')\\
\frac{\cos\varphi}{\cos h}\sin(t''-k)&=\sin(A+A'-A'')
\end{aligned}
\]}
so gilt\HistoricalNote{6}
\[
\varphi=D+E,\qquad h=D-E
\]

Ein Nachteil dieser Methode ist, dass sich schwer im Voraus eine bestimmte Regel dafür angeben lässt, in welchem Halbkreis man die Winkel $A,A',A'',B,D,E$ nehmen muss, wenn man allein nach den analytischen Formeln ohne Zeichnung rechnet und den Beobachtungen nur die Uhrzeiten entnimmt, also vorher nicht weiß, auf welchen Seiten des Meridians die drei Sterne beobachtet wurden. Bekanntlich lassen die Tangenten allein diese Wahl offen. Es lässt sich jedoch zeigen, dass man zunächst beliebig wählen darf. In einigen Fällen muss man dann noch die folgenden Änderungen vornehmen.

1) Man ersetze $k$ durch $k+180^0$ oder, was hier auf dasselbe hinausläuft, durch $k-180^0$, wenn die für $\varphi$ und $h$ gefundenen Werte dazu führen, dass $\cos\varphi$ und $\sin h$ entgegengesetzte Vorzeichen haben.
\SourceLeaf{149}%

2) Liegen die für $h$ und $\varphi$ gefundenen Werte außerhalb der Grenzen $0$ und $90^0$, so nimmt man ihre Differenz zum nächstgelegenen Vielfachen von $180^0$.

3) Die Polhöhe ist als nördlich oder südlich anzusehen, je nachdem, ob $\sin\varphi$ und $\sin h$ gleiche oder entgegengesetzte Vorzeichen haben.

Ergäben also zum Beispiel die obigen Formeln $\varphi=231^0$, $h=-127^0$, so wäre die tatsächliche Polhöhe $51^0$ und die tatsächliche Höhe $53^0$. Der Wert von $k$ bliebe dabei gleich. Hätte man dagegen $\varphi=231^0$, $h=+127^0$ gefunden, so müsste man $k$ um $180^0$ vergrößern oder verkleinern. Das sei natürlich nur der analytischen Vollständigkeit halber bemerkt. Wenn $k$ um $180^0$ geändert werden muss, ist das ohnehin klar, da man beim Uhrstand nicht um ganze 12 Stunden im Zweifel ist.

Zum Vergleich lege ich noch die Rechnung mit den obigen Formeln für dasselbe Zahlenbeispiel bei.

\clearpage\begin{center}\emph{Berechnung des Beispiels nach Cagnolis Methode.}\end{center}
\begin{center}\small
\begin{minipage}[t]{.49\textwidth}\centering
\begin{tabular}[t]{@{}lr@{}}
$\tfrac12(\delta''-\delta')=$&$-24^0\ 49'\ 59''55$\\
$\tfrac12(\delta-\delta'')=$&$-5\quad17\quad25.90$\\
$\tfrac12(\delta'-\delta)=$&$+30\quad7\quad25.45$\\[10pt]
$\tfrac12(t''-t')=$&$-129\quad41\quad39.45$\\
$\tfrac12(t-t'')=$&$+135\quad0\quad4.72$\\
$\tfrac12(t'-t)=$&$-5\quad18\quad25.27$
\end{tabular}
\end{minipage}\hfill
\begin{minipage}[t]{.49\textwidth}\centering
\begin{tabular}[t]{@{}lr@{}}
$\tfrac12(\delta''+\delta')=$&$63^0\ 27'\ 6''15$\\
$\tfrac12(\delta+\delta'')=$&$33\quad19\quad40.70$\\
$\tfrac12(\delta'+\delta)=$&$58\quad9\quad40.25$
\end{tabular}
\end{minipage}
\par\medskip
\begin{tabular}[b]{@{}lr@{}}
$\log\sin\tfrac12(\delta''-\delta')\;.\;.\;.$&$9.6232267\,\mathrm n$\\
$\mathrm{Compl.}\log\cos\tfrac12(\delta''+\delta')\;.$&$0.3497391$\\
$\log\operatorname{cotang}\tfrac12(t''-t)\;.\;.\;.$&$9.9191030$\\\hline
$\log\operatorname{tang}A\;.\;.\;.\;.\;.\;.$&$9.8920688\,\mathrm n$
\end{tabular}
\qquad $A=142^0\quad2\quad50''70$
\par\medskip
\begin{tabular}[b]{@{}lr@{}}
$\log\sin\tfrac12(\delta-\delta'')\;.\;.\;.\;.$&$8.9647590\,\mathrm n$\\
$\mathrm{Compl.}\log\cos\tfrac12(\delta+\delta'')\;.$&$0.0780332$\\
$\log\operatorname{cotang}\tfrac12(t-t')\;.\;.\;.$&$0.0000199\,\mathrm n$\\\hline
$\log\operatorname{tang}A'\;.\;.\;.\;.\;.$&$9.0428121$
\end{tabular}
\qquad $A'=6^0\ 17'\ 51''34$
\end{center}
\SourceLeaf{150}%
\begin{center}\begingroup\small\setlength{\tabcolsep}{2.5pt}
\sbox{0}{\begin{tabular}{@{}lr@{\hspace{16pt}}rl@{}}
$\log\sin\tfrac12(\delta'-\delta)\;.\;.\;.\;.$&$9.7005892$&&\\
$\mathrm{Compl.}\log\cos\tfrac12(\delta'+\delta)\;.$&$0.2777516$&&\\\cline{1-2}
$(*)\;.\;.\;.\;.\;.\;.\;.\;.$&$9.9783408$&&\\
$\log\operatorname{cotang}\tfrac12(t'-t)\;.\;.\;.$&$1.0320274$&&\\\cline{1-2}
$\log\operatorname{tang}A''\;.\;.\;.\;.\;.$&$1.0103682\,\mathrm n$&$A''=$&$95^0\ 34'\ 36''24$\\
&&$A-A'=$&$135\quad44\quad59.36$\\
&&$C=$&$-40\quad10\quad23.12$\\[7pt]
$(*)\;.\;.\;.\;.\;.\;.\;.\;.$&$9.9783408$&&\\
$\log\operatorname{cotang}(A-A')\;.\;.\;.$&$0.0113684\,\mathrm n$&&\\\cline{1-2}
$\log\operatorname{tang}B\;.\;.\;.\;.\;.$&$9.9897092\,\mathrm n$&$B=$&$315^0\ 40'\ 43''55$\\
&&$\tfrac12(t+t')=$&$318\quad24\quad44.78$\\\cline{3-4}
&&$k=$&$2\quad44\quad1.23$\\[4pt]
&&$\tfrac12(t-k+C)=$&$140\quad24\quad22.85$\\
&&$\tfrac12(t-k-C)=$&$180\quad34\quad45.97$\\
&&$45^0+\tfrac12\delta=$&$59\quad1\quad7.40$\\[7pt]
$\log\cos\tfrac12(t-k+C)\;.\;.$&$9.8868199\,\mathrm n$&&\\
$\mathrm{Compl.}\log\cos\tfrac12(t-k-C)$&$0.0000222\,\mathrm n$&&\\
$\log\operatorname{tang}(45^0+\tfrac12\delta)\;.\;.\;.$&$0.2215478$&&\\\cline{1-2}
$\log\operatorname{tang}D\;.\;.\;.\;.\;.$&$0.1083899$&$D=$&$52^0\ 4'\ 36''35$\\[7pt]
$\log\sin\tfrac12(t-k-C)\;.\;.$&$8.0048736\,\mathrm n$&&\\
$\mathrm{Compl.}\log\sin\tfrac12(t-k+C)$&$0.1956297$&&\\
$\log\operatorname{cotang}(45^0+\tfrac12\delta)\;.\;.$&$9.7784522\,\mathrm n$&&\\\cline{1-2}
$\log\operatorname{tang}E\;.\;.\;.\;.\;.$&$7.9789555\,\mathrm n$&$E=$&$-0^0\ 32'\ 45''02$\\
&&Also $\varphi=$&$51\quad31\quad51.33$\\
&&$h=$&$52\quad37\quad21.37$\\
\end{tabular}}
\ifdim\wd0>\linewidth\resizebox{\linewidth}{!}{\usebox{0}}\else\usebox{0}\fi
\endgroup\end{center}
\par\bigskip\centerline{\rule{.27\linewidth}{.4pt}}
